\documentclass[14pt,a4paper]{extarticle}
\usepackage[T2A]{fontenc}
\usepackage[utf8]{inputenc}
\usepackage[russian]{babel}
\usepackage{amsmath, amssymb, amsfonts}
\usepackage{graphicx}
\usepackage{geometry}
\usepackage{titlesec}
\usepackage{fancyhdr}
\usepackage{xcolor}
\usepackage{tikz}
\usepackage{pgfplots}
\usepackage{minted}
\usepackage{caption}
\usepackage{float}
\usepackage{hyperref}

\usemintedstyle{vs}
\setminted{
    linenos=true,
    fontsize=\scriptsize,
    frame=single,
    framesep=2mm,
    breaklines=true,
    breakanywhere=true,
    tabsize=4,
    bgcolor=gray!10
}
\begin{document}
\begin{titlepage}
\begin{center}
МИНИСТЕРСТВО НАУКИ И ВЫСШЕГО ОБРАЗОВАНИЯ\\
РОССИЙСКОЙ ФЕДЕРАЦИИ\\
Федеральное государственное бюджетное образовательное учреждение\\
высшего образования\\[0.5cm]
{\Large\bfseries АДЫГЕЙСКИЙ ГОСУДАРСТВЕННЫЙ УНИВЕРСИТЕТ}\\[0.3cm]
НОК ``Институт точных наук и цифровых технологий''\\
Кафедра автоматизированных систем обработки информации и управления\\
\vfill
{\Large ОТЧЕТ ПО ПРАКТИЧЕСКОЙ РАБОТЕ}\\[0.5cm]
{\Large Программная реализация численного метода}\\
{\Large\itshape Вариант 7. Нахождение определителя матрицы}\\[0.5cm]
1 курс, группа 1ИВТ2\\
\vfill
\begin{flushright}
\begin{tabular}{l}
Выполнил: \\
\hrulefill\ Ю.\,И.\ Цирульник \\
«\underline{\hspace{1cm}}» \underline{\hspace{2cm}} 2026 г. \\[0.5cm]
Проверил: \\
\hrulefill\ \underline{\hspace{6cm}} \\
«\underline{\hspace{1cm}}» \underline{\hspace{2cm}} 2026 г. \\
\end{tabular}
\end{flushright}
\vfill
Майкоп, 2026 г.
\end{center}
\end{titlepage}


\section{Теоретические сведения}

\subsection*{Постановка задачи}
Необходимо разработать программу для вычисления определителя квадратной матрицы произвольного порядка $n$.

\subsection*{Теория}
Определителем квадратной матрицы $A = \begin{pmatrix} a_{11} & a_{12} \\ a_{21} & a_{22} \end{pmatrix}$ второго порядка называется число:
\begin{equation}
|A| = a_{11}a_{22} - a_{12}a_{21}.
\label{eq:det2}
\end{equation}

Определителем $A = \begin{pmatrix} a_{11} & \cdots & a_{1n} \\ \vdots & \ddots & \vdots \\ a_{n1} & \cdots & a_{nn} \end{pmatrix}$ квадратной матрицы порядка $n$, $n \geq 3$, называется число:
\begin{equation}
|A| = \sum_{k=1}^{n} (-1)^{k+1} a_{1k} M_k,
\label{eq:detn}
\end{equation}
где $M_k$ --- определитель матрицы порядка $n-1$, полученной из матрицы $A$ вычёркиванием первой строки и столбца с номером $k$ (минор элемента $a_{1k}$).

\section{Ход выполнения работы}

\subsection{Разработка программы}
Программа реализована на языке Python. Она включает в себя:
\begin{itemize}
    \item Модуль ввода данных: Включает функции, которые обеспечивают пошаговый ввод размерности матрицы и её элементов. Реализована проверка корректности вводимых данных (проверка на целое число для размера, на числовой формат элементов и на соответствие количества введенных значений размеру строки).
    \item Алгоритм вычисления: Используется рекурсивная функция, которая вычисляет определитель методом разложения по первой строке согласно формуле. Для формирования миноров применяется вспомогательная функция.
    \item Обработка ошибок: Программа обрабатывает прерывание выполнения пользователем (KeyboardInterrupt) и перехватывает непредвиденные исключения, обеспечивая стабильную работу приложения.
\end{itemize}

\subsection{Исходный код}
Ниже представлен исходный код программы \texttt{main.py}:

\begin{minted}[caption={Исходный код программы вычисления определителя}]{python}
def read_matrix_size():
    while True:
        try:
            n = int(input("Введите порядок квадратной матрицы (N): "))
            if n > 0:
                return n
            print("Ошибка: размерность должна быть положительным числом.")
        except ValueError:
            print("Ошибка: введите целое число.")


def read_matrix_row(size, row_index):
    while True:
        try:
            row = list(map(float, input(f"Строка {row_index}: ").split()))
            if len(row) != size:
                print(f"Ошибка: ожидается {size} чисел, введено {len(row)}.")
                continue
            return row
        except ValueError:
            print("Ошибка: допустимы только числа.")


def build_square_matrix(n):
    print(f"Введите элементы матрицы {n}x{n} (числа через пробел):")
    return [read_matrix_row(n, i + 1) for i in range(n)]


def get_submatrix(matrix, remove_row, remove_col):
    return [
        [value for col_idx, value in enumerate(row) if col_idx != remove_col]
        for row_idx, row in enumerate(matrix) if row_idx != remove_row
    ]


def calculate_determinant(matrix):
    n = len(matrix)

    if n == 1:
        return matrix[0][0]
    if n == 2:
        return matrix[0][0] * matrix[1][1] - matrix[0][1] * matrix[1][0]

    det_value = 0
    for col_idx in range(n):
        sign = 1 if col_idx % 2 == 0 else -1
        element = matrix[0][col_idx]
        minor = get_submatrix(matrix, 0, col_idx)

        det_value += sign * element * calculate_determinant(minor)

    return det_value

def format_result(value):
    return int(value) if value.is_integer() else value

def main():
    try:
        size = read_matrix_size()
        matrix = build_square_matrix(size)

        result = calculate_determinant(matrix)
        print(f"Определитель равен: {format_result(result)}")

    except KeyboardInterrupt:
        print("Программа прервана пользователем.")
    except Exception as e:
        print(f"Неожиданная ошибка: {e}")

if __name__ == "__main__":
    main()
\end{minted}

\section{Результаты работы}

В данном разделе приведены примеры работы программы для матриц различной размерности.

\begin{figure}[H]
\centering
\begin{figure}
    \centering
    \includegraphics[width=1\linewidth]{matrix_2x2.png}
    \caption{Результат выполнения матрицы 2х2.}
    \label{fig:placeholder}
\end{figure}
\label{fig:res2}
\end{figure}

\begin{figure}[H]
\centering
\begin{figure}
    \centering
    \includegraphics[width=1\linewidth]{matrix_3x3.png}
    \caption{Результат выполнения матрицы 3х3.}
    \label{fig:placeholder}
\end{figure}
\label{fig:res3}
\end{figure}

\begin{figure}[H]
\centering
\begin{figure}
    \centering
    \includegraphics[width=1\linewidth]{matrix_4x4.png}
    \caption{Результат выполнения матрицы 4х4.}
    \label{fig:placeholder}
\end{figure}
\label{fig:res4}
\end{figure}

Как видно из рисунков 1-3, программа корректно вычисляет определитель для матриц различных порядков и обрабатывает некорректный ввод пользователя, что соответствует требованиям задания.

\section{Выводы}
В ходе работы была реализована программа на языке Python для нахождения определителя матрицы. Использовался метод разложения по первой строке. Реализован контроль вводимых данных, что повышает универсальность программы.

\clearpage

\begin{thebibliography}{9}

\bibitem{ilin}
Ильин В.А., Позняк Э.Г. Линейная алгебра. --- 3-е изд., перераб. и доп. --- Москва: Наука, 1984. --- 296 с.

\bibitem{strang}
Стрэнг Г. Линейная алгебра и её приложения. --- Санкт-Петербург: Лань, 2013. --- 464 с.

\bibitem{python_doc}
Документация языка Python 3.12. \url{https://docs.python.org/3/} (дата обращения: 20.05.2024).

\bibitem{svetlova}
Светлова Е.Н. Python. Быстрый старт. --- Москва: ДМК Пресс, 2018. --- 272 с.

\end{thebibliography}

\end{document}